\documentclass[11pt]{article}

\usepackage[margin=1in]{geometry}
\usepackage{amsmath,amssymb,mathtools}
\usepackage{bm}
\usepackage{booktabs}
\usepackage{microtype}
\usepackage[colorlinks=true,linkcolor=blue]{hyperref}

\newcommand{\ii}{\mathrm{i}}
\newcommand{\ee}{\mathrm{e}}
\newcommand{\cQ}{\mathcal{Q}}
\newcommand{\cC}{\mathcal{C}}
\newcommand{\cL}{\mathcal{L}}
\newcommand{\cD}{\mathcal{D}}
\newcommand{\cR}{\mathcal{R}}
\newcommand{\ip}[2]{\left\langle #1,#2\right\rangle}
\newcommand{\cc}{\mathrm{c.c.}}
\newcommand{\T}{\mathbb{T}}
\newcommand{\one}[1]{\bm{1}_{\{#1\}}}
\DeclareMathOperator{\RePart}{Re}
\DeclareMathOperator{\spanop}{span}

\title{Weakly Nonlinear Reduction with Separated Temporal and Spatial Forcing}
\author{}
\date{\today}

\begin{document}

\maketitle

\section{Problem setting}

We consider
\begin{align}
 u_t={}&\cL(p)u+\cQ(u,u)+\cC(u,u,u) \notag\\
 &+\epsilon^n\left(P\ee^{\ii\omega t}+\overline P\ee^{-\ii\omega t}\right)
 +\epsilon^m\left(S\ee^{\ii kx}+\overline S\ee^{-\ii kx}\right),
 \label{eq:pde}
\end{align}
where $u(x,t)\in\mathbb R^d$, $x\in\T=\mathbb R/(2\pi\mathbb Z)$,
$p\in\mathbb R$, and
\begin{equation}
 \cL(p)=\cL_0+p\cL_1+O(p^2).
 \label{eq:Lexpansion}
\end{equation}
The forcing vectors are denoted by $P,S$ so that $\cQ$ is reserved for the
quadratic map.

Assume:
\begin{enumerate}
 \item $\cQ$ and $\cC$ are symmetric multilinear maps, extended complex
 linearly;
 \item $\cL_j$ is real, translation invariant, and reflection invariant,
 $j=0,1$;
 \item the only critical modes are the simple finite-wavenumber Hopf modes
 $(\Omega,h)=(\pm\omega_0,\pm k_0)$, where $\omega_0>0$ and
 $k_0\in\mathbb N$.
\end{enumerate}
Here reflection invariance is the intended symmetry of $\cL_j$; it does not
mean that $\cL_j$ is self-adjoint.

For a spatial Fourier mode, define
\begin{equation}
 \cL_j\left(v\ee^{\ii hx}\right)=M_j(h)v\ee^{\ii hx},
 \qquad h\in\mathbb Z.
 \label{eq:Fourier-symbol}
\end{equation}
Reality and reflection symmetry give, respectively,
\begin{equation}
 M_j(-h)=\overline{M_j(h)},
 \qquad M_j(-h)=M_j(h).
 \label{eq:M-symmetries}
\end{equation}
Thus $M_j(h)$ is real in this setting and the matrices at $h$ and $-h$ are
identical.

Choose right and adjoint eigenvectors $\Phi,\Psi\in\mathbb C^d$ such that
\begin{equation}
 M_0(\pm k_0)\Phi=\ii\omega_0\Phi,
 \qquad
 M_0(\pm k_0)^*\Psi=-\ii\omega_0\Psi,
 \qquad
 \Psi^*\Phi=1.
 \label{eq:eigenvectors}
\end{equation}
The conjugate equations give the eigenvectors $\overline\Phi$ and
$\overline\Psi$ at eigenvalues $-\ii\omega_0$ and $+\ii\omega_0$ of
$M_0$ and $M_0^*$, respectively.  Because of
\eqref{eq:M-symmetries}, the same $\Phi,\Psi$ can be used at $\pm k_0$.

Let $\nu(h,p)$ be the eigenvalue branch through
$\nu(\pm k_0,0)=\ii\omega_0$.  If
$M(\pm k_0,p)\Phi(p)=\nu(\pm k_0,p)\Phi(p)$, differentiation at $p=0$
gives
\begin{equation}
 M_1(\pm k_0)\Phi+M_0(\pm k_0)\Phi_p
 =\nu_p\Phi+\ii\omega_0\Phi_p.
 \label{eq:eigenvalue-differentiation}
\end{equation}
Since $\Psi^*M_0(\pm k_0)=\ii\omega_0\Psi^*$, multiplication by $\Psi^*$
cancels both terms containing $\Phi_p$.  Therefore
\begin{equation}
 \partial_p\nu(\pm k_0,0)
 =\Psi^*M_1(\pm k_0)\Phi=: \alpha,
 \qquad \RePart\alpha>0.
 \label{eq:transversality}
\end{equation}
The last inequality is the transverse-crossing assumption.  In particular,
$\alpha\ne0$.

Define the modal resolvent
\begin{equation}
 R(\Omega,h):=\left(\ii\Omega I-M_0(h)\right)^{-1}.
 \label{eq:resolvent}
\end{equation}
Because the underlying system is real,
$M_0(-h)=\overline{M_0(h)}$, and hence
\begin{equation}
 R(-\Omega,-h)=\overline{R(\Omega,h)}.
 \label{eq:resolvent-conjugacy}
\end{equation}
In particular,
\begin{equation}
 R(-\omega,0)=\overline{R(\omega,0)},
 \qquad
 R(0,-k)=\overline{R(0,k)}.
 \label{eq:forcing-resolvent-conjugacy}
\end{equation}
With the reflection symmetry assumed here, $M_0(-k)=M_0(k)$ as well, so
$R(0,-k)=R(0,k)=\overline{R(0,k)}$.
The no-quadratic-resonance assumption is
\begin{equation}
 \begin{gathered}
 2\ii\omega_0I-M_0(0),\quad
 2\ii\omega_0I-M_0(\pm2k_0),\\
 -M_0(0),\quad -M_0(\pm2k_0)
 \end{gathered}
 \quad\text{are invertible.}
 \label{eq:no-quadratic-resonance}
\end{equation}
We additionally assume that $R(\pm\omega,0)$ and $R(0,\pm k)$ exist whenever
the corresponding forcing is present.  Failure of any of these
invertibilities introduces an additional center mode and requires a larger
amplitude system.

Introduce fast and slow times
\begin{equation}
 s=t,
 \qquad \tau=\epsilon^2t,
 \qquad \partial_t=\partial_s+\epsilon^2\partial_\tau,
 \label{eq:slow-time}
\end{equation}
and initially write
\begin{equation}
 p=\epsilon p_1+\epsilon^2p_2+O(\epsilon^3),
 \qquad
 u=\epsilon u_1+\epsilon^2u_2+\epsilon^3u_3+O(\epsilon^4).
 \label{eq:expansions}
\end{equation}
The terms that determine the first three equations are
\begin{align}
 \partial_tu={}&\epsilon\partial_su_1+\epsilon^2\partial_su_2
 +\epsilon^3(\partial_su_3+\partial_\tau u_1)+O(\epsilon^4),
 \label{eq:expand-time}\\
 \cL_0u={}&\epsilon\cL_0u_1+\epsilon^2\cL_0u_2
 +\epsilon^3\cL_0u_3+O(\epsilon^4),
 \label{eq:expand-L0}\\
 p\cL_1u={}&\epsilon^2p_1\cL_1u_1
 +\epsilon^3(p_2\cL_1u_1+p_1\cL_1u_2)+O(\epsilon^4),
 \label{eq:expand-L1}\\
 \cQ(u,u)={}&\epsilon^2\cQ(u_1,u_1)
 +2\epsilon^3\cQ(u_1,u_2)+O(\epsilon^4),
 \label{eq:expand-Q}\\
 \cC(u,u,u)={}&\epsilon^3\cC(u_1,u_1,u_1)+O(\epsilon^4).
 \label{eq:expand-C}
\end{align}
The factor $2$ in \eqref{eq:expand-Q} is the symmetry of $\cQ$.  The term
$p_1\cL_1u_2$ and every higher-$p$ contribution at $O(\epsilon^3)$ vanish
after the second-order solvability condition gives $p_1=0$.

\section{The solution space of
\texorpdfstring{$\cD=\partial_s-\cL_0$}{D = ds - L0}}

Set
\begin{equation}
 \cD:=\partial_s-\cL_0,
 \qquad
 E_\pm:=\ee^{\ii(\omega_0s\pm k_0x)}.
 \label{eq:D-and-E}
\end{equation}
For any pure wave,
\begin{equation}
 \cD\left(v\ee^{\ii(\Omega s+hx)}\right)
 =\left(\ii\Omega I-M_0(h)\right)v
 \ee^{\ii(\Omega s+hx)}.
 \label{eq:D-on-mode}
\end{equation}
The multiplier in \eqref{eq:D-on-mode} is singular at
$(\Omega,h)=(\omega_0,\pm k_0)$ with nullvector $\Phi$, and its conjugate is
singular at $(-\omega_0,\mp k_0)$ with nullvector $\overline\Phi$.  There are
no other singular Fourier blocks by assumption.  Hence
\begin{equation}
 \ker\cD
 =\spanop_{\mathbb R}\left\{
 \Phi E_+,\Phi E_-,\overline\Phi\,\overline E_+,
 \overline\Phi\,\overline E_-
 \right\}.
 \label{eq:kernel-D}
\end{equation}

On functions periodic in $x$ and in $s$ with
$T_0=2\pi/\omega_0$, use
\begin{equation}
 \ip{v}{w}:=\frac{1}{2\pi T_0}
 \int_0^{T_0}\int_0^{2\pi}v(s,x)^*w(s,x)\,dx\,ds.
 \label{eq:inner-product}
\end{equation}
Periodic integration by parts gives
\begin{equation}
 \cD^*=-\partial_s-\cL_0^*.
 \label{eq:D-star}
\end{equation}
Indeed,
\begin{equation}
 \cD^*\left(w\ee^{\ii(\Omega s+hx)}\right)
 =\left(-\ii\Omega I-M_0(h)^*\right)w
 \ee^{\ii(\Omega s+hx)},
 \label{eq:Dstar-on-mode}
\end{equation}
so \eqref{eq:eigenvectors} gives
\begin{equation}
 \ker\cD^*=\spanop_{\mathbb R}\left\{
 \Psi E_+,\Psi E_-,\overline\Psi\,\overline E_+,
 \overline\Psi\,\overline E_-
 \right\}.
 \label{eq:D-adjoint}
\end{equation}

For pure waves,
\begin{equation}
 \ip{a\ee^{\ii(\Omega s+hx)}}
 {b\ee^{\ii(\Omega' s+h'x)}}
 =a^*b\,\delta_{\Omega,\Omega'}\delta_{h,h'}.
 \label{eq:wave-orthogonality}
\end{equation}
Thus different wave types are orthogonal.  Having the same
frequency--wavenumber pair is necessary, but not sufficient, for a nonzero
inner product because $a^*b$ may also vanish.  For incommensurate forcing
frequencies, \eqref{eq:wave-orthogonality} is interpreted as a long-time mean;
equivalently, one works on the torus of independent fast phases.

Let
\begin{equation}
 U_\pm=\Phi E_\pm,
 \qquad V_\pm=\Psi E_\pm.
\end{equation}
The complete pairings needed below are
\begin{equation}
\begin{array}{c|cccc}
 &\cL_1U_+&\cL_1U_-&\cL_1\overline U_+&\cL_1\overline U_-\\ \hline
 V_+&\alpha&0&0&0\\
 V_-&0&\alpha&0&0\\
 \overline V_+&0&0&\overline\alpha&0\\
 \overline V_-&0&0&0&\overline\alpha
\end{array}.
 \label{eq:L1-pairing-table}
\end{equation}

Under the spectral assumptions above, $\cD$ is Fredholm with
finite-dimensional kernel and closed range: in Fourier variables it is the
direct sum of the blocks in \eqref{eq:D-on-mode}, all noncritical blocks are
invertible, and the critical part is finite dimensional.  Hence
\begin{equation}
 \operatorname{Ran}\cD=(\ker\cD^*)^\perp,
 \qquad
 \cD w=f\ \Longleftrightarrow\
 \ip{V}{f}=0\quad\text{for every }V\in\ker\cD^*.
 \label{eq:fredholm}
\end{equation}
At every order, the goal is to put the right-hand side in
$\operatorname{Ran}\cD$.

To prevent the stimulus itself from supplying the leading pattern, take
$n,m>1$.  If $n=1$ or $m=1$, the corresponding stimulus appears in
$\cD u_1$ and $u_1$ is no longer chosen solely from $\ker\cD$.  If an
exponent equals $2$, its response is part of $u_2$ and interacts with $u_1$
through $2\cQ(u_1,u_2)$ at order three.  An exponent greater than $2$ cannot
produce such a cubic-order feedback.  In particular, exponent-$3$ forcing
would enter \eqref{eq:order3} directly, but its separated types $(\omega,0)$
and $(0,\pm k)$ are never critical because $k_0,\omega_0>0$; it changes only
$u_3$.  Exponents larger than $3$ do not occur in the present truncation.
Thus only exponent-$2$ forcing can modify the leading amplitude equations,
and we define
\begin{align}
 F_2={}&\one{n=2}
 \left(P\ee^{\ii\omega s}+\overline P\ee^{-\ii\omega s}\right)
 \notag\\
 &+\one{m=2}
 \left(S\ee^{\ii kx}+\overline S\ee^{-\ii kx}\right).
 \label{eq:F2}
\end{align}
The three equations are
\begin{subequations}
\label{eq:three-orders}
\begin{align}
 O(\epsilon):\quad
 &\cD u_1=0,
 \label{eq:order1}\\
 O(\epsilon^2):\quad
 &\cD u_2=p_1\cL_1u_1+\cQ(u_1,u_1)+F_2,
 \label{eq:order2}\\
 O(\epsilon^3):\quad
 &\cD u_3=-\partial_\tau u_1+p_2\cL_1u_1
 +2\cQ(u_1,u_2)+\cC(u_1,u_1,u_1),
 \label{eq:order3}
\end{align}
\end{subequations}
where \eqref{eq:order3} is written after the second-order condition below
sets $p_1=0$.  We subsequently write $p_2=\mu$.

\section{\texorpdfstring{$O(\epsilon)$}{O(epsilon)} solution}

The real leading solution is
\begin{equation}
 \boxed{
 u_1=A(\tau)\Phi E_+ +B(\tau)\Phi E_-+\cc.}
 \label{eq:u1}
\end{equation}
The amplitudes $A$ and $B$ are independent.  The states $B=0$ and $A=0$ are
the two traveling-wave directions.  Equal moduli $|A|=|B|$ give standing
waves; imposing $A=B$ at the outset would restrict the calculation to one
standing-wave subspace.

\section{\texorpdfstring{$O(\epsilon^2)$}{O(epsilon squared)} solution}

The right-hand side of \eqref{eq:order2} must be orthogonal to
$\ker\cD^*$.  Adding any two wave types from
$(\pm\omega_0,\pm k_0)$ produces one of the types in
\eqref{eq:quadratic-wave-types} below, never $(\omega_0,\pm k_0)$.  Hence
\begin{equation}
 \ip{V_\pm}{\cQ(u_1,u_1)}=0.
 \label{eq:Q-second-projection}
\end{equation}
Neither separated forcing is a critical wave because
$(\omega,0)\ne(\omega_0,\pm k_0)$ and
$(0,\pm k)\ne(\omega_0,\pm k_0)$, so
\begin{equation}
 \ip{V_\pm}{F_2}=0.
 \label{eq:F2-projection}
\end{equation}
Only $p_1\cL_1u_1$ remains.  By
\eqref{eq:L1-pairing-table},
\begin{align}
 \ip{V_+}{p_1\cL_1u_1}&=p_1A\alpha,
 &\ip{V_-}{p_1\cL_1u_1}&=p_1B\alpha.
 \label{eq:p1-projections}
\end{align}
Thus the two independent complex conditions are
\begin{equation}
 0=p_1A\alpha,
 \qquad 0=p_1B\alpha.
 \label{eq:p1-conditions}
\end{equation}
For $(A,B)\ne(0,0)$ and $\RePart\alpha>0$,
\begin{equation}
 \boxed{p_1=0,\qquad p=\epsilon^2\mu+O(\epsilon^3).}
 \label{eq:p1-zero}
\end{equation}

Define
\begin{equation}
 q_{20}:=\cQ(\Phi,\Phi),
 \qquad
 q_{11}:=\cQ(\Phi,\overline\Phi).
 \label{eq:q-vectors}
\end{equation}
To track the factors of two, set
\begin{equation}
 u_1^+:=\Phi(AE_++BE_-),
 \qquad u_1=u_1^++\overline{u_1^+}.
\end{equation}
Symmetry of $\cQ$ gives
\begin{equation}
 \cQ(u_1,u_1)=\cQ(u_1^+,u_1^+)
 +2\cQ(u_1^+,\overline{u_1^+})
 +\cQ(\overline{u_1^+},\overline{u_1^+}),
 \label{eq:Q-three-parts}
\end{equation}
and therefore
\begin{align}
 \cQ(u_1,u_1)={}&q_{20}
 \left(A^2E_+^2+B^2E_-^2+2AB E_+E_-\right)+\cc
 \notag\\
 &+2q_{11}\left(
 |A|^2+|B|^2+A\overline B E_+\overline E_-
 +\overline A B\overline E_+E_-
 \right).
 \label{eq:Q-u1-u1}
\end{align}
The generated wave types are
\begin{equation}
 (\pm2\omega_0,\pm2k_0),\quad
 (\pm2\omega_0,0),\quad
 (0,\pm2k_0),\quad(0,0).
 \label{eq:quadratic-wave-types}
\end{equation}
They are noncritical by wave-type arithmetic; condition
\eqref{eq:no-quadratic-resonance} is the additional statement that $\cD$ is
invertible on each of them.

For a noncritical type, \eqref{eq:D-on-mode} gives the mode-by-mode rule
\begin{equation}
 \cD^{-1}\left(f\ee^{\ii(\Omega s+hx)}\right)
 =R(\Omega,h)f\ee^{\ii(\Omega s+hx)}.
 \label{eq:D-inverse-rule}
\end{equation}

Set
\begin{subequations}
\label{eq:h-vectors}
\begin{align}
 h_{20}^+&=R(2\omega_0,2k_0)q_{20},
 &h_{20}^-&=R(2\omega_0,-2k_0)q_{20},\\
 h_{20}^0&=2R(2\omega_0,0)q_{20},
 &h_{11}^0&=2R(0,0)q_{11},\\
 h_{11}^+&=2R(0,2k_0)q_{11},
 &h_{11}^-&=2R(0,-2k_0)q_{11}.
\end{align}
\end{subequations}
The six nonconjugate source-to-response steps are
\begin{equation}
\begin{array}{c|c|c}
 \text{source in }\cQ(u_1,u_1)&(\Omega,h)&\text{term in }u_2^Q\\ \hline
 A^2q_{20}E_+^2&(2\omega_0,2k_0)&A^2h_{20}^+E_+^2\\
 B^2q_{20}E_-^2&(2\omega_0,-2k_0)&B^2h_{20}^-E_-^2\\
 2ABq_{20}E_+E_-&(2\omega_0,0)&ABh_{20}^0E_+E_-\\
 2(|A|^2+|B|^2)q_{11}&(0,0)&(|A|^2+|B|^2)h_{11}^0\\
 2A\overline Bq_{11}E_+\overline E_-&(0,2k_0)
 &A\overline Bh_{11}^+E_+\overline E_-\\
 2\overline ABq_{11}\overline E_+E_-&(0,-2k_0)
 &\overline ABh_{11}^-\overline E_+E_-
\end{array}.
 \label{eq:u2-source-response-table}
\end{equation}
Then
\begin{align}
 u_2^{Q}:={}&\cD^{-1}\cQ(u_1,u_1)\notag\\
 ={}&A^2h_{20}^+E_+^2+B^2h_{20}^-E_-^2
 +ABh_{20}^0E_+E_-+\cc\notag\\
 &+(|A|^2+|B|^2)h_{11}^0
 +A\overline B h_{11}^+E_+\overline E_-
 +\overline A B h_{11}^-\overline E_+E_-.
 \label{eq:u2Q}
\end{align}
Reflection symmetry gives $h_{20}^+=h_{20}^-$ and
$h_{11}^+=h_{11}^-$.  The superscripts are retained to display the wave
types.

For the general forcing frequency and wavenumber, define the
positive-frequency and positive-wavenumber response coefficients
\begin{equation}
\label{eq:forcing-responses}
 p_\omega:=R(\omega,0)P,
 \qquad
 s_k:=R(0,k)S.
\end{equation}
The symmetry \eqref{eq:forcing-resolvent-conjugacy} gives
\begin{equation}
 R(-\omega,0)\overline P=\overline{p_\omega},
 \qquad
 R(0,-k)\overline S=\overline{s_k}.
 \label{eq:forcing-response-conjugacy}
\end{equation}
Thus the negative-frequency temporal coefficient and the negative-wavenumber
spatial coefficient are the complex conjugates of the corresponding positive
coefficients; they are not independent quantities.
For example,
\begin{equation}
 \cD\left(p_\omega\ee^{\ii\omega s}\right)
 =(\ii\omega I-M_0(0))p_\omega\ee^{\ii\omega s}
 =P\ee^{\ii\omega s},
\end{equation}
and the other three forcing modes are identical applications of
\eqref{eq:D-inverse-rule}.  Thus
\begin{align}
 u_2^F={}&\cD^{-1}F_2\notag\\
 ={}&\one{n=2}\left(
 p_\omega\ee^{\ii\omega s}+\overline{p_\omega}\ee^{-\ii\omega s}\right)
 +\one{m=2}\left(
 s_k\ee^{\ii kx}+\overline{s_k}\ee^{-\ii kx}\right).
 \label{eq:u2F}
\end{align}
The conjugate terms make $u_2^F$ real.  We use the normalization
\begin{equation}
 \boxed{u_2=u_2^Q+u_2^F,}
 \label{eq:u2-total}
\end{equation}
with no added component from $\ker\cD$.

\section{\texorpdfstring{$O(\epsilon^3)$}{O(epsilon cubed)} solution and
solvability}

With $p_2=\mu$, write the right-hand side of \eqref{eq:order3} in four
groups:
\begin{align}
 F_{0,\mathrm{lin}}^{(3)}
 &:=-\partial_\tau u_1+\mu\cL_1u_1,
 \label{eq:F0lin}\\
 F_{0,\mathrm{int}}^{(3)}
 &:=2\cQ(u_1,u_2^Q)+\cC(u_1,u_1,u_1),
 \label{eq:F0int}\\
 F_{F,\mathrm{lin}}^{(3)}&:=0,
 \label{eq:FFlin}\\
 F_{F,\mathrm{int}}^{(3)}&:=2\cQ(u_1,u_2^F).
 \label{eq:FFint}
\end{align}
The third line is zero because both retained stimuli entered at order two;
their only cubic-order effect is through $u_2^F$.

\subsection{Unforced terms}

For the critical wave $E_+$,
\begin{equation}
 [F_{0,\mathrm{lin}}^{(3)}]_{E_+}
 =-A_\tau\Phi+\mu A M_1(k_0)\Phi.
 \label{eq:unforced-linear-plus}
\end{equation}
The resonant cubic contribution is
\begin{equation}
 [\cC(u_1,u_1,u_1)]_{E_+}
 =3A|A|^2\cC(\Phi,\Phi,\overline\Phi)
 +6A|B|^2\cC(\Phi,\Phi,\overline\Phi).
 \label{eq:C-resonant}
\end{equation}
Indeed, a triple of critical phases can sum to $E_+$ only as
$E_+E_+\overline E_+$ or $E_+E_-\overline E_-$.  Their permutation counts
are $3$ and $6$, respectively, which gives \eqref{eq:C-resonant}.

For the quadratic feedback, frequency and wavenumber must add separately.
The complete $E_+$ selection table is
\begin{equation}
\begin{array}{c|c|c}
 \text{factor from }u_1&\text{factor from }u_2^Q
 &\text{coefficient after }2\cQ\\ \hline
 A\Phi E_+&(|A|^2+|B|^2)h_{11}^0
 &2A(|A|^2+|B|^2)\cQ(\Phi,h_{11}^0)\\
 \overline A\,\overline\Phi\,\overline E_+&A^2h_{20}^+E_+^2
 &2A|A|^2\cQ(\overline\Phi,h_{20}^+)\\
 \overline B\,\overline\Phi\,\overline E_-&ABh_{20}^0E_+E_-
 &2A|B|^2\cQ(\overline\Phi,h_{20}^0)\\
 B\Phi E_-&A\overline Bh_{11}^+E_+\overline E_-
 &2A|B|^2\cQ(\Phi,h_{11}^+)
\end{array}.
 \label{eq:quadratic-Eplus-selection}
\end{equation}
No other pair of wave types sums to $(\omega_0,k_0)$.  Summing the last
column gives
\begin{align}
 [2\cQ(u_1,u_2^Q)]_{E_+}={}&
 2A(|A|^2+|B|^2)\cQ(\Phi,h_{11}^0)
 \notag\\
 &+2A|A|^2\cQ(\overline\Phi,h_{20}^+)
 +2A|B|^2\cQ(\overline\Phi,h_{20}^0)
 \notag\\
 &+2A|B|^2\cQ(\Phi,h_{11}^+).
 \label{eq:Q-feedback}
\end{align}
Define
\begin{subequations}
\label{eq:G-vectors}
\begin{align}
 G_{\mathrm{s}}={}&2\cQ(\Phi,h_{11}^0)
 +2\cQ(\overline\Phi,h_{20}^+)
 +3\cC(\Phi,\Phi,\overline\Phi),
 \label{eq:Gs}\\
 G_{\mathrm{c}}={}&2\cQ(\Phi,h_{11}^0)
 +2\cQ(\overline\Phi,h_{20}^0)
 +2\cQ(\Phi,h_{11}^+)
 +6\cC(\Phi,\Phi,\overline\Phi).
 \label{eq:Gc}
\end{align}
\end{subequations}
Therefore
\begin{subequations}
\label{eq:unforced-critical-coefficients}
\begin{align}
 [F_{0,\mathrm{lin}}^{(3)}+F_{0,\mathrm{int}}^{(3)}]_{E_+}
 ={}&-A_\tau\Phi+\mu A M_1(k_0)\Phi
 +A|A|^2G_{\mathrm{s}}+A|B|^2G_{\mathrm{c}},
 \label{eq:Eplus-coeff}\\
 [F_{0,\mathrm{lin}}^{(3)}+F_{0,\mathrm{int}}^{(3)}]_{E_-}
 ={}&-B_\tau\Phi+\mu B M_1(k_0)\Phi
 +B|B|^2G_{\mathrm{s}}+B|A|^2G_{\mathrm{c}}.
 \label{eq:Eminus-coeff}
\end{align}
\end{subequations}
These are the only unforced terms outside $\operatorname{Ran}\cD$.

\subsection{Forcing-interaction terms and their wave types}

Let $a_+=A$, $a_-=B$, and let $\sigma\in\{+1,-1\}$.  The complete temporal
forcing interaction is
\begin{align}
 2\cQ(u_1,u_{2,t}^F)=2\one{n=2}\sum_{\sigma=\pm1}\Big[{}
 &a_\sigma\cQ(\Phi,p_\omega)
 \ee^{\ii((\omega_0+\omega)s+\sigma k_0x)}
 \notag\\
 +{}&a_\sigma\cQ(\Phi,\overline{p_\omega})
 \ee^{\ii((\omega_0-\omega)s+\sigma k_0x)}
 \notag\\
 +{}&\overline{a_\sigma}\cQ(\overline\Phi,p_\omega)
 \ee^{\ii((\omega-\omega_0)s-\sigma k_0x)}
 \notag\\
 +{}&\overline{a_\sigma}\cQ(\overline\Phi,\overline{p_\omega})
 \ee^{-\ii((\omega+\omega_0)s+\sigma k_0x)}
 \Big].
 \label{eq:temporal-interaction-full}
\end{align}
Temporal forcing does not change spatial wavenumber.  Therefore an $E_+$
term can only use the two first-order factors whose wavenumber is $+k_0$:
\begin{equation}
 \begin{array}{c|c|c}
 \text{first-order type}&\text{matching equation}&\text{solution}\\ \hline
 (\omega_0,k_0)&(\omega_0,k_0)+(\pm\omega,0)
 =(\omega_0,k_0)&\omega=0\\
 (-\omega_0,k_0)&(-\omega_0,k_0)+(\omega,0)
 =(\omega_0,k_0)&\omega=2\omega_0
 \end{array}.
 \label{eq:temporal-selection-logic}
\end{equation}
The $E_-$ calculation is its reflection.
For $\omega\ge0$, its positive-frequency critical pieces are
\begin{align}
 [2\cQ(u_1,u_{2,t}^F)]_{E_+}
 ={}&\one{n=2}\left[
 \one{\omega=0}\,2A\cQ(\Phi,p_0+\overline{p_0})
 +\one{\omega=2\omega_0}\,2\overline B
 \cQ(\overline\Phi,p_{2\omega_0})\right],
 \label{eq:temporal-critical-A}\\
 [2\cQ(u_1,u_{2,t}^F)]_{E_-}
 ={}&\one{n=2}\left[
 \one{\omega=0}\,2B\cQ(\Phi,p_0+\overline{p_0})
 +\one{\omega=2\omega_0}\,2\overline A
 \cQ(\overline\Phi,p_{2\omega_0})\right].
 \label{eq:temporal-critical-B}
\end{align}
The forcing terms that survive Fredholm projection therefore have fast wave
types
\begin{equation}
 \boxed{
 (2\omega_0,0)+(-\omega_0,\pm k_0)
 =(\omega_0,\pm k_0).}
 \label{eq:temporal-wave-match}
\end{equation}

The complete spatial forcing interaction is
\begin{align}
 2\cQ(u_1,u_{2,x}^F)=2\one{m=2}\sum_{\sigma=\pm1}\Big[{}
 &a_\sigma\cQ(\Phi,s_k)
 \ee^{\ii(\omega_0s+(\sigma k_0+k)x)}
 \notag\\
 +{}&a_\sigma\cQ(\Phi,\overline{s_k})
 \ee^{\ii(\omega_0s+(\sigma k_0-k)x)}
 \notag\\
 +{}&\overline{a_\sigma}\cQ(\overline\Phi,s_k)
 \ee^{\ii(-\omega_0s+(-\sigma k_0+k)x)}
 \notag\\
 +{}&\overline{a_\sigma}\cQ(\overline\Phi,\overline{s_k})
 \ee^{-\ii(\omega_0s+(\sigma k_0+k)x)}
 \Big].
 \label{eq:spatial-interaction-full}
\end{align}
Spatial forcing does not change temporal frequency.  Hence an $E_+$ term
must use a positive-frequency first-order factor, and the possible matches
are
\begin{equation}
 \begin{array}{c|c|c}
 \text{first-order type}&\text{matching equation}&\text{solution}\\ \hline
 (\omega_0,k_0)&(\omega_0,k_0)+(0,\pm k)
 =(\omega_0,k_0)&k=0\\
 (\omega_0,-k_0)&(\omega_0,-k_0)+(0,k)
 =(\omega_0,k_0)&k=2k_0
 \end{array}.
 \label{eq:spatial-selection-logic}
\end{equation}
Again, reflection gives the $E_-$ calculation.
For $k\ge0$, its positive-frequency critical pieces are
\begin{align}
 [2\cQ(u_1,u_{2,x}^F)]_{E_+}
 ={}&\one{m=2}\left[
 \one{k=0}\,2A\cQ(\Phi,s_0+\overline{s_0})
 +\one{k=2k_0}\,2B\cQ(\Phi,s_{2k_0})\right],
 \label{eq:spatial-critical-A}\\
 [2\cQ(u_1,u_{2,x}^F)]_{E_-}
 ={}&\one{m=2}\left[
 \one{k=0}\,2B\cQ(\Phi,s_0+\overline{s_0})
 +\one{k=2k_0}\,2A\cQ(\Phi,\overline{s_{2k_0}})\right].
 \label{eq:spatial-critical-B}
\end{align}
The nontrivial spatial match is
\begin{equation}
 \boxed{
 (0,\pm2k_0)+(\omega_0,\mp k_0)
 =(\omega_0,\pm k_0).}
 \label{eq:spatial-wave-match}
\end{equation}

Consequently, relative to the unforced problem, the complete list of
cubic-order exceptions for $n,m>1$ is
\begin{equation}
 \begin{array}{c|c|c}
 \text{forcing}&\text{condition}&\text{critical term}\\ \hline
 \text{temporal DC}&n=2,\ \omega=0&A,\ B\\
 \text{temporal parametric}&n=2,\ \omega=2\omega_0&\overline B,\ \overline A\\
 \text{spatial DC}&m=2,\ k=0&A,\ B\\
 \text{spatial parametric}&m=2,\ k=2k_0&B,\ A
 \end{array}.
 \label{eq:forcing-case-table}
\end{equation}
If $\omega$ and $k$ are restricted to be strictly positive, the two DC rows
are absent.
For all other $\omega,k$, every term in
\eqref{eq:temporal-interaction-full} and
\eqref{eq:spatial-interaction-full} lies in $\operatorname{Ran}\cD$ and only
changes $u_3$.

\subsection{Fredholm conditions}

Define
\begin{equation}
 a:=\mu\Psi^*M_1(k_0)\Phi=\mu\alpha,
 \qquad b:=\Psi^*G_{\mathrm{s}},
 \qquad c:=\Psi^*G_{\mathrm{c}}.
 \label{eq:abc}
\end{equation}
Because of the normalized fast phases and $\Psi^*\Phi=1$,
\begin{equation}
 \ip{V_+}{F^{(3)}}=\Psi^*[F^{(3)}]_{E_+},
 \qquad
 \ip{V_-}{F^{(3)}}=\Psi^*[F^{(3)}]_{E_-}.
 \label{eq:projection-extracts-coefficient}
\end{equation}
For completeness, the forcing projections are
\begin{subequations}
\label{eq:forcing-projections}
\begin{align}
 \mathcal T_A={}&\one{n=2}\left[
 \one{\omega=0}\,2A\Psi^*\cQ(\Phi,p_0+\overline{p_0})
 +\one{\omega=2\omega_0}\,2\overline B
 \Psi^*\cQ(\overline\Phi,p_{2\omega_0})\right],\\
 \mathcal T_B={}&\one{n=2}\left[
 \one{\omega=0}\,2B\Psi^*\cQ(\Phi,p_0+\overline{p_0})
 +\one{\omega=2\omega_0}\,2\overline A
 \Psi^*\cQ(\overline\Phi,p_{2\omega_0})\right],\\
 \mathcal X_A={}&\one{m=2}\left[
 \one{k=0}\,2A\Psi^*\cQ(\Phi,s_0+\overline{s_0})
 +\one{k=2k_0}\,2B\Psi^*\cQ(\Phi,s_{2k_0})\right],\\
 \mathcal X_B={}&\one{m=2}\left[
 \one{k=0}\,2B\Psi^*\cQ(\Phi,s_0+\overline{s_0})
 +\one{k=2k_0}\,2A\Psi^*\cQ(\Phi,\overline{s_{2k_0}})\right].
\end{align}
\end{subequations}
Projection onto $V_+$ and $V_-$ now gives
\begin{subequations}
\label{eq:full-solvability}
\begin{align}
 0={}&-A_\tau+A(a+b|A|^2+c|B|^2)+\mathcal T_A+\mathcal X_A,
 \label{eq:solvability-A}\\
 0={}&-B_\tau+B(a+b|B|^2+c|A|^2)+\mathcal T_B+\mathcal X_B,
 \label{eq:solvability-B}
\end{align}
\end{subequations}
The conjugate adjoint modes give only the conjugates of
\eqref{eq:full-solvability}; hence these are the complete independent
conditions.

\section{Amplitude equations and pattern selection}

\subsection{No leading forcing resonance}

When the forcing projections vanish, the amplitude equations are
\begin{subequations}
\label{eq:unforced-amplitude}
\begin{align}
 A_\tau&=A(a+b|A|^2+c|B|^2),
 \label{eq:A-unforced}\\
 B_\tau&=B(a+b|B|^2+c|A|^2).
 \label{eq:B-unforced}
\end{align}
\end{subequations}
Write $a_r=\RePart a$, $b_r=\RePart b$, and $c_r=\RePart c$.  For
$A=r_+\ee^{\ii\theta_+}$ and $B=r_-\ee^{\ii\theta_-}$,
\begin{subequations}
\label{eq:radial-system}
\begin{align}
 (r_+)_\tau&=r_+(a_r+b_rr_+^2+c_rr_-^2),\\
 (r_-)_\tau&=r_-(a_r+b_rr_-^2+c_rr_+^2).
\end{align}
\end{subequations}
The zero state has the double radial eigenvalue $a_r$ and is stable for
$a_r<0$.  For $a_r>0$, a traveling wave satisfies
\begin{equation}
 r_- =0,
 \qquad a_r+b_rr_+^2=0,
 \qquad \rho^2:=r_+^2=-\frac{a_r}{b_r}.
 \label{eq:TW-balance}
\end{equation}
Its radial and transverse eigenvalues are
\begin{equation}
 \lambda_{\mathrm{rad}}=-2a_r,
 \qquad
 \lambda_{\mathrm{tr}}=a_r+c_r\rho^2
 =a_r\frac{b_r-c_r}{b_r}.
 \label{eq:TW-eigenvalues}
\end{equation}
A standing wave satisfies
\begin{equation}
 r_+=r_-=\rho,
 \qquad
 \rho^2=-\frac{a_r}{b_r+c_r}.
 \label{eq:SW-balance}
\end{equation}
The radial Jacobian is
\begin{equation}
 2\rho^2
 \begin{pmatrix}b_r&c_r\\c_r&b_r\end{pmatrix},
 \label{eq:SW-Jacobian}
\end{equation}
with symmetric and antisymmetric eigenvalues
\begin{equation}
 \lambda_{\mathrm{sym}}=-2a_r,
 \qquad
 \lambda_{\mathrm{asym}}=2\rho^2(b_r-c_r).
 \label{eq:SW-eigenvalues}
\end{equation}
Consequently the nonzero equilibria modulo phase are
\begin{equation}
\begin{array}{c|c|c|c}
 \text{state}&(r_+,r_-)&\text{existence}&\text{amplitude stability}\\ \hline
 \text{traveling wave}&(\rho,0)\text{ or }(0,\rho)
 &\rho^2=-a_r/b_r>0&c_r-b_r<0\\
 \text{standing wave}&(\rho,\rho)
 &\rho^2=-a_r/(b_r+c_r)>0&c_r-b_r>0
\end{array}.
 \label{eq:state-table}
\end{equation}
The imaginary parts give the phase equations
\begin{align}
 (\theta_+)_\tau&=\Im(a)+\Im(b)r_+^2+\Im(c)r_-^2,\\
 (\theta_-)_\tau&=\Im(a)+\Im(b)r_-^2+\Im(c)r_+^2,
\end{align}
and hence only phase drift and frequency corrections in the unforced
problem.  The leading physical patterns are
\begin{align}
 u_{\mathrm{TW}}(x,t)
 &\sim2\epsilon\rho\RePart\left[
 \Phi\ee^{\ii(\omega_0t\pm k_0x+\theta)}\right],
 \label{eq:TW-pattern}\\
 u_{\mathrm{SW}}(x,t)
 &\sim4\epsilon\rho\cos(k_0x+\vartheta_x)
 \RePart\left[
 \Phi\ee^{\ii(\omega_0t+\vartheta_t)}\right].
 \label{eq:SW-pattern}
\end{align}

\subsection{Leading forcing resonances}

The DC responses merely shift the linear coefficient.  Define
\begin{align}
 \chi_0={}&2\one{n=2}\one{\omega=0}\,
 \Psi^*\cQ(\Phi,p_0+\overline{p_0})
 \notag\\
 &+2\one{m=2}\one{k=0}\,
 \Psi^*\cQ(\Phi,s_0+\overline{s_0}).
 \label{eq:DC-shift}
\end{align}
Then $a$ is replaced by $a+\chi_0$ in both equations.

For temporal $2{:}1$ forcing, take
\begin{equation}
 n=2,
 \qquad \omega=2\omega_0+\epsilon^2\delta_t,
 \qquad
 \eta_t:=2\Psi^*\cQ\left(
 \overline\Phi,R(2\omega_0,0)P\right).
 \label{eq:temporal-coefficient}
\end{equation}
The slow factor is $\ee^{\ii\delta_t\tau}$, and
\begin{subequations}
\label{eq:temporal-amplitude}
\begin{align}
 A_\tau&=A(a+b|A|^2+c|B|^2)
 +\eta_t\ee^{\ii\delta_t\tau}\overline B,\\
 B_\tau&=B(a+b|B|^2+c|A|^2)
 +\eta_t\ee^{\ii\delta_t\tau}\overline A.
\end{align}
\end{subequations}

For spatial $2{:}1$ forcing, take
\begin{equation}
 m=2,
 \qquad k=2k_0,
 \qquad
 \eta_x^+=2\Psi^*\cQ(\Phi,R(0,2k_0)S),
 \quad
 \eta_x^-=2\Psi^*\cQ(\Phi,R(0,-2k_0)\overline S).
 \label{eq:spatial-coefficients}
\end{equation}
Then
\begin{subequations}
\label{eq:spatial-amplitude}
\begin{align}
 A_\tau&=A(a+b|A|^2+c|B|^2)+\eta_x^+B,\\
 B_\tau&=B(a+b|B|^2+c|A|^2)+\eta_x^-A.
\end{align}
\end{subequations}

If both $2{:}1$ conditions hold, the contributions add:
\begin{subequations}
\label{eq:combined-amplitude}
\begin{align}
 A_\tau={}&A(a+\chi_0+b|A|^2+c|B|^2)
 +\eta_t\ee^{\ii\delta_t\tau}\overline B+\eta_x^+B,
 \label{eq:combined-A}\\
 B_\tau={}&B(a+\chi_0+b|B|^2+c|A|^2)
 +\eta_t\ee^{\ii\delta_t\tau}\overline A+\eta_x^-A.
 \label{eq:combined-B}
\end{align}
\end{subequations}
With $A=\widetilde A\ee^{\ii\delta_t\tau/2}$ and
$B=\widetilde B\ee^{\ii\delta_t\tau/2}$, this becomes autonomous after the
replacement $a+\chi_0\mapsto a+\chi_0-\ii\delta_t/2$.  Explicitly,
\begin{equation}
 A_\tau=\ee^{\ii\delta_t\tau/2}
 \left(\widetilde A_\tau+\frac{\ii\delta_t}{2}\widetilde A\right),
 \qquad
 \ee^{\ii\delta_t\tau}\overline B
 =\ee^{\ii\delta_t\tau/2}\overline{\widetilde B},
\end{equation}
and similarly for the second equation.

The decomposition \eqref{eq:u2F} is linear, so at cubic order
\begin{equation}
 2\cQ(u_1,u_2^F)
 =2\cQ(u_1,u_{2,t}^F)+2\cQ(u_1,u_{2,x}^F).
 \label{eq:forcing-additivity}
\end{equation}
The first temporal--spatial product occurs one order later:
\begin{equation}
 2\cQ(\epsilon^2u_{2,t}^F,\epsilon^2u_{2,x}^F)
 =O(\epsilon^4).
 \label{eq:no-forcing-cross-term}
\end{equation}
Thus \eqref{eq:combined-amplitude} contains the complete leading combination.
Temporal forcing couples a mode to the conjugate counter-propagating mode;
spatial forcing couples the two nonconjugate traveling modes.  Their
phase-locked asymptotic states are the roots of the right-hand side of the
autonomous form of \eqref{eq:combined-amplitude}; stability is determined by
its $4\times4$ real Jacobian.  No additional universal stability inequality
exists without the complex values of $a,b,c,\eta_t,\eta_x^\pm$.

\end{document}
